\chapter*{Conclusion}
\addcontentsline{toc}{chapter}{Conclusion}
\markboth{Conclusion}{Conclusion}

Ce stage opérateur au sein de Reminex, filiale d'ingénierie du Groupe Managem, m'a permis d'atteindre l'objectif que je m'étais fixé~: découvrir, de l'intérieur, le monde industriel et le métier d'ingénieur.

J'y ai compris la \textbf{complexité d'un projet minier}, qui mobilise une longue chaîne d'étapes, de l'exploration jusqu'à la production, et un grand nombre de disciplines qui doivent travailler ensemble. J'ai découvert ce qu'est concrètement l'\textbf{ingénierie de procédé} --- « transformer un \textit{flowsheet} en usine » --- ainsi que le langage documentaire et la logique économique qui la sous-tendent. J'ai surtout mesuré l'\textbf{importance de la collaboration} et de la communication entre les métiers, sans lesquelles aucun projet ne peut aboutir.

Sur le plan personnel, cette expérience a fait évoluer ma vision de l'ingénierie, désormais indissociable, à mes yeux, de sa dimension humaine et organisationnelle. Elle a renforcé mon intérêt pour l'ingénierie industrielle et la gestion de projet, et elle éclaire le choix de mes prochains stages, que je souhaite plus proches de la conception et de la conduite de projet.

Je retiens enfin une leçon plus large~: on apprend beaucoup en observant, à condition de le faire activement, en questionnant et en cherchant à relier ce que l'on voit à une vision d'ensemble. C'est dans cet état d'esprit que j'aborderai la suite de ma formation à l'École Centrale Casablanca.
